\chapter{附录 C：A++ 如何对应CMMI-DEV ML4 要求} % Introduction chapter suppressed from the table of contents

\hypertarget{ux56e2ux961fux7ea7}{%
\subsection{团队级}\label{ux56e2ux961fux7ea7}}

\hypertarget{est-2.1-ux5236ux5b9aux4f7fux7528ux5e76ux4fddux6301ux66f4ux65b0ux4f30ux7b97ux8303ux56f4}{%
\subsubsection{EST 2.1
制定、使用并保持更新估算范围}\label{est-2.1-ux5236ux5b9aux4f7fux7528ux5e76ux4fddux6301ux66f4ux65b0ux4f30ux7b97ux8303ux56f4}}

\begin{itemize}
\tightlist
\item
  每轮迭代，在大白纸上写上每人2-4周每天的活动安排。
\item
  从大白纸桑的活动安排，看到本迭代的WBS，每个工作包（活动）的开始、结束日期和工作量（人天）。{[}7
  章{]}
\end{itemize}

\hypertarget{est-2.2-ux5236ux5b9aux5e76ux4fddux6301ux66f4ux65b0ux9488ux5bf9ux89e3ux51b3ux65b9ux6848ux89c4ux6a21ux7684ux4f30ux7b97}{%
\subsubsection{EST 2.2
制定并保持更新针对解决方案规模的估算}\label{est-2.2-ux5236ux5b9aux5e76ux4fddux6301ux66f4ux65b0ux9488ux5bf9ux89e3ux51b3ux65b9ux6848ux89c4ux6a21ux7684ux4f30ux7b97}}

\begin{itemize}
\tightlist
\item
  估算规模

  \begin{itemize}
  \tightlist
  \item
    使用简化功能点估算开发和维护项目规模。{[}第六章{]}
  \end{itemize}
\end{itemize}

\hypertarget{est-2.3-ux6839ux636eux89c4ux6a21ux4f30ux7b97ux6765ux5236ux5b9aux5e76ux8bb0ux5f55ux89e3ux51b3ux65b9ux6848ux6240ux9700ux5de5ux4f5cux91cfux5468ux671fux548cux6210ux672cux53caux5176ux4f9dux636e}{%
\subsubsection{EST 2.3
根据规模估算来制定并记录解决方案所需工作量、周期和成本及其依据}\label{est-2.3-ux6839ux636eux89c4ux6a21ux4f30ux7b97ux6765ux5236ux5b9aux5e76ux8bb0ux5f55ux89e3ux51b3ux65b9ux6848ux6240ux9700ux5de5ux4f5cux91cfux5468ux671fux548cux6210ux672cux53caux5176ux4f9dux636e}}

\hypertarget{est-3.1-ux5236ux5b9aux5e76ux4fddux6301ux66f4ux65b0ux5df2ux8bb0ux5f55ux7684ux4f30ux7b97ux65b9ux6cd5}{%
\subsubsection{EST 3.1
制定并保持更新已记录的估算方法}\label{est-3.1-ux5236ux5b9aux5e76ux4fddux6301ux66f4ux65b0ux5df2ux8bb0ux5f55ux7684ux4f30ux7b97ux65b9ux6cd5}}

\begin{itemize}
\tightlist
\item
  制定估算过程和指南

  \begin{itemize}
  \tightlist
  \item
    参考国际功能点、简化功能点的指南和实例。{[}第六章{]}
  \end{itemize}
\end{itemize}

\hypertarget{est-3.2-ux4f7fux7528ux7ec4ux7ec7ux7684ux5ea6ux91cfux5e93ux548cux8fc7ux7a0bux8d44ux4ea7ux8fdbux884cux4f30ux7b97ux5de5ux4f5c}{%
\subsubsection{EST 3.2
使用组织的度量库和过程资产进行估算工作}\label{est-3.2-ux4f7fux7528ux7ec4ux7ec7ux7684ux5ea6ux91cfux5e93ux548cux8fc7ux7a0bux8d44ux4ea7ux8fdbux884cux4f30ux7b97ux5de5ux4f5c}}

\begin{itemize}
\tightlist
\item
  依据项目历史数据估算工作量

  \begin{itemize}
  \tightlist
  \item
    项目组收集每个迭代的总工作量和规模大小得出生产力画成趋势图。如果稳定下来，就可以参考作为项目的标杆，在以后的项目能更准确帮助工作量估算。
  \end{itemize}
\end{itemize}

\hypertarget{plan-2.1-ux5f00ux53d1ux5b8cux6210ux5de5ux4f5cux7684ux65b9ux6cd5ux5e76ux4fddux6301ux66f4ux65b0}{%
\subsubsection{PLAN 2.1
开发完成工作的方法并保持更新}\label{plan-2.1-ux5f00ux53d1ux5b8cux6210ux5de5ux4f5cux7684ux65b9ux6cd5ux5e76ux4fddux6301ux66f4ux65b0}}

\begin{itemize}
\tightlist
\item
  项目生命周期

  \begin{itemize}
  \tightlist
  \item
    按需求稳定性、团队能力、对进度/工作量/质量等项目要求，选择合适的项目生命周期。{[}7
    章{]}
  \end{itemize}
\end{itemize}

\hypertarget{plan-2.3-ux6839ux636eux8bb0ux5f55ux7684ux4f30ux7b97ux5236ux5b9aux9884ux7b97ux548cux8fdbux5ea6ux5e76ux4fddux6301ux66f4ux65b0}{%
\subsubsection{PLAN 2.3
根据记录的估算，制定预算和进度并保持更新}\label{plan-2.3-ux6839ux636eux8bb0ux5f55ux7684ux4f30ux7b97ux5236ux5b9aux9884ux7b97ux548cux8fdbux5ea6ux5e76ux4fddux6301ux66f4ux65b0}}

\begin{itemize}
\tightlist
\item
  策划项目预算

  \begin{itemize}
  \tightlist
  \item
    如果公司关注项目成本，尤其是专门为甲方做定制开发的团队，往往在项目结束后才知道项目超支，如果要管理项目的成本，就必须要使用WBS项目管理工具，在立项的时候输进每个里程碑的预算是多少。然后工具按每周员工填工时表和每个工种的月薪，计算实际成本，用挣值分析监控成本偏差指数。{[}第七章{]}
  \end{itemize}
\end{itemize}

工具也可以帮助项目监控成本和进度偏差更满足PLAN2.6，人力资源安排到活动2.7和3.3
活动之间的依赖关系和MC2.1实际与计划对比。

项目如果是走迭代的话，虽然迭代里面可以用大白纸简单估算。但整个项目就还需要使用项目管理工具才可以及时得出偏差指数

\hypertarget{plan-2.6-ux6839ux636eux8d44ux6e90ux7684ux5bb9ux91cfux548cux53efux7528ux6027ux534fux8c03ux4f30ux7b97ux786eux4fddux8ba1ux5212ux7684ux53efux884cux6027}{%
\subsubsection{PLAN 2.6
根据资源的容量和可用性协调估算，确保计划的可行性}\label{plan-2.6-ux6839ux636eux8d44ux6e90ux7684ux5bb9ux91cfux548cux53efux7528ux6027ux534fux8c03ux4f30ux7b97ux786eux4fddux8ba1ux5212ux7684ux53efux884cux6027}}

\begin{itemize}
\tightlist
\item
  分配人手到活动

  \begin{itemize}
  \tightlist
  \item
    每轮迭代，在大白纸上写上每人2-4周每天的活动安排。{[}7 章{]}
  \item
    从大白纸桑的活动安排，看到本迭代的WBS，每个工作包（活动）的开始、结束日期和工作量（人天）。{[}7
    章{]}
  \end{itemize}
\end{itemize}

\hypertarget{plan-2.7-ux5236ux5b9aux9879ux76eeux8ba1ux5212ux786eux4fddux5176ux8981ux7d20ux4e4bux95f4ux7684ux4e00ux81f4ux6027ux5e76ux4fddux6301ux66f4ux65b0}{%
\subsubsection{PLAN 2.7
制定项目计划，确保其要素之间的一致性，并保持更新}\label{plan-2.7-ux5236ux5b9aux9879ux76eeux8ba1ux5212ux786eux4fddux5176ux8981ux7d20ux4e4bux95f4ux7684ux4e00ux81f4ux6027ux5e76ux4fddux6301ux66f4ux65b0}}

\hypertarget{plan-3.3-ux8bc6ux522bux548cux534fux5546ux5173ux952eux4f9dux8d56ux5173ux7cfb}{%
\subsubsection{PLAN 3.3
识别和协商关键依赖关系}\label{plan-3.3-ux8bc6ux522bux548cux534fux5546ux5173ux952eux4f9dux8d56ux5173ux7cfb}}

\begin{itemize}
\tightlist
\item
  管理活动间依赖关系，确保子计划一致

  \begin{itemize}
  \tightlist
  \item
    每轮迭代，用大白纸便利贴管理活动间依赖关系。也能确保开发与测试活动之间没有冲突。
  \end{itemize}
\end{itemize}

\hypertarget{plan-4.1-ux4f7fux7528ux7edfux8ba1ux4e0eux5176ux4ed6ux91cfux5316ux6280ux672fux6765ux5f00ux53d1ux9879ux76eeux8fc7ux7a0bux5e76ux4fddux6301ux66f4ux65b0ux4ee5ux5b9eux73b0ux8d28ux91cfux4e0eux8fc7ux7a0bux6027ux80fdux76eeux6807}{%
\subsubsection{PLAN 4.1
使用统计与其他量化技术来开发项目过程并保持更新,以实现质量与过程性能目标}\label{plan-4.1-ux4f7fux7528ux7edfux8ba1ux4e0eux5176ux4ed6ux91cfux5316ux6280ux672fux6765ux5f00ux53d1ux9879ux76eeux8fc7ux7a0bux5e76ux4fddux6301ux66f4ux65b0ux4ee5ux5b9eux73b0ux8d28ux91cfux4e0eux8fc7ux7a0bux6027ux80fdux76eeux6807}}

\begin{itemize}
\tightlist
\item
  制定用于评价项目过程备选方案的准则

  \begin{itemize}
  \tightlist
  \item
    总成本最低。
  \end{itemize}
\item
  识别或开发可用来完成工作和实现目标的替代过程

  \begin{itemize}
  \tightlist
  \item
    模型中已进行体现，通过不同的评审方法互相替换。
  \end{itemize}
\item
  根据记录的评价标准分析和评价备选过程

  \begin{itemize}
  \tightlist
  \item
    参考蒙特卡洛Optimize 仿真选择最佳配搭。
  \end{itemize}
\item
  选择最符合准则的备选过程

  \begin{itemize}
  \tightlist
  \item
    同上，在模型中体现。
  \end{itemize}
\item
  评价无法实现项目的质量与过程性能目标的风险

  \begin{itemize}
  \tightlist
  \item
    利用PPM预测结果的分布，分析概率，
  \end{itemize}
\end{itemize}

\hypertarget{mc-2.1-ux4eceux89c4ux6a21ux5de5ux4f5cux91cfux8fdbux5ea6ux8d44ux6e90ux77e5ux8bc6ux548cux6280ux80fdux4ee5ux53caux9884ux7b97ux7b49ux65b9ux9762ux5bf9ux6bd4ux4f30ux7b97ux8ddfux8e2aux5b9eux9645ux7ed3ux679c}{%
\subsubsection{MC 2.1
从规模、工作量、进度、资源、知识和技能以及预算等方面,对比估算跟踪实际结果}\label{mc-2.1-ux4eceux89c4ux6a21ux5de5ux4f5cux91cfux8fdbux5ea6ux8d44ux6e90ux77e5ux8bc6ux548cux6280ux80fdux4ee5ux53caux9884ux7b97ux7b49ux65b9ux9762ux5bf9ux6bd4ux4f30ux7b97ux8ddfux8e2aux5b9eux9645ux7ed3ux679c}}

\begin{itemize}
\tightlist
\item
  对比实际与计划，如进度、工作量、缺陷数等

  \begin{itemize}
  \tightlist
  \item
    2-4周迭代：在大白纸上写上实际
  \end{itemize}
\end{itemize}

\hypertarget{mc-2.4-ux5f53ux5b9eux9645ux7ed3ux679cux76f8ux8f83ux4e8eux8ba1ux5212ux7ed3ux679cux5b58ux5728ux663eux8457ux5deeux5f02ux65f6ux91c7ux53d6ux7ea0ux6b63ux63aaux65bdux5e76ux7ba1ux7406ux76f4ux81f3ux5173ux95ed}{%
\subsubsection{MC 2.4
当实际结果相较于计划结果存在显著差异时，采取纠正措施并管理直至关闭}\label{mc-2.4-ux5f53ux5b9eux9645ux7ed3ux679cux76f8ux8f83ux4e8eux8ba1ux5212ux7ed3ux679cux5b58ux5728ux663eux8457ux5deeux5f02ux65f6ux91c7ux53d6ux7ea0ux6b63ux63aaux65bdux5e76ux7ba1ux7406ux76f4ux81f3ux5173ux95ed}}

\begin{itemize}
\tightlist
\item
  分析项目的问题，根本原因，采取纠正措施。

  \begin{itemize}
  \tightlist
  \item
    如果项目团队在每个迭代回顾时，针对缺陷排除率和返工做根因分析，并在下个迭代采取纠正措施，就是这条的实践。{[}第十章-第十三章{]}
  \end{itemize}
\end{itemize}

\hypertarget{mc-3.1-ux4f7fux7528ux9879ux76eeux8ba1ux5212ux548cux9879ux76eeux8fc7ux7a0bux7ba1ux7406ux9879ux76ee}{%
\subsubsection{MC 3.1
使用项目计划和项目过程管理项目}\label{mc-3.1-ux4f7fux7528ux9879ux76eeux8ba1ux5212ux548cux9879ux76eeux8fc7ux7a0bux7ba1ux7406ux9879ux76ee}}

\begin{itemize}
\tightlist
\item
  按实际活动完成情况更新计划

  \begin{itemize}
  \tightlist
  \item
    在大白纸上写上活动实际完成日期，让干系人，包括团队成员，知道本迭代进度偏差。
  \item
    每轮迭代后， 计算当前累计的SPI，CPI，也写在大白纸上。
  \end{itemize}
\end{itemize}

\hypertarget{mpm-3.1-ux5f00ux53d1ux6301ux7eedux66f4ux65b0ux548cux4f7fux7528ux53efux8ffdux6eafux5230ux4e1aux52a1ux76eeux6807ux7684ux7ec4ux7ec7ux5ea6ux91cfux53caux6027ux80fdux76eeux6807}{%
\subsubsection{MPM 3.1
开发、持续更新和使用可追溯到业务目标的组织度量及性能目标}\label{mpm-3.1-ux5f00ux53d1ux6301ux7eedux66f4ux65b0ux548cux4f7fux7528ux53efux8ffdux6eafux5230ux4e1aux52a1ux76eeux6807ux7684ux7ec4ux7ec7ux5ea6ux91cfux53caux6027ux80fdux76eeux6807}}

\hypertarget{mpm-4.1-ux4f7fux7528ux7edfux8ba1ux4e0eux5176ux4ed6ux91cfux5316ux6280ux672fux6765ux5236ux5b9aux6301ux7eedux66f4ux65b0ux5e76ux6c9fux901aux53efux8ffdux6eafux5230ux4e1aux52a1ux76eeux6807ux7684ux8d28ux91cfux4e0eux8fc7ux7a0bux6027ux80fdux76eeux6807}{%
\subsubsection{MPM 4.1
使用统计与其他量化技术来制定，持续更新并沟通可追溯到业务目标的质量与过程性能目标}\label{mpm-4.1-ux4f7fux7528ux7edfux8ba1ux4e0eux5176ux4ed6ux91cfux5316ux6280ux672fux6765ux5236ux5b9aux6301ux7eedux66f4ux65b0ux5e76ux6c9fux901aux53efux8ffdux6eafux5230ux4e1aux52a1ux76eeux6807ux7684ux8d28ux91cfux4e0eux8fc7ux7a0bux6027ux80fdux76eeux6807}}

\begin{itemize}
\tightlist
\item
  质量过程目标

  \begin{itemize}
  \tightlist
  \item
    组织: 用控制图、蒙特卡罗模型、方差分析等来分析SMART质量过程目标
  \item
    项目: 用预测模型、基线为下个冲刺预测质量过程目标
  \end{itemize}
\item
  制定中期目标以监控进度

  \begin{itemize}
  \tightlist
  \item
    项目:
    用预测模型来预测下一个冲刺的需求审查，代码审查，单元测试的缺陷密度范围
  \end{itemize}
\item
  确定并记录风险

  \begin{itemize}
  \tightlist
  \item
    使用预测模型来预测实现质量过程目标的百分比
  \end{itemize}
\item
  解决质量过程目标之间的冲突

  \begin{itemize}
  \tightlist
  \item
    预测模型使用整体成本来优化，可以平衡质量（减少化系统测试/验收测试缺陷密度）和生产率目标之间的冲突
  \end{itemize}
\end{itemize}

\hypertarget{mpm-4.3-ux4f7fux7528ux7edfux8ba1ux4e0eux5176ux4ed6ux91cfux5316ux6280ux672fux6765ux5efaux7acbux548cux5206ux6790ux8fc7ux7a0bux6027ux80fdux57faux7ebfux5e76ux4fddux6301ux66f4ux65b0}{%
\subsubsection{MPM 4.3
使用统计与其他量化技术来建立和分析过程性能基线并保持更新}\label{mpm-4.3-ux4f7fux7528ux7edfux8ba1ux4e0eux5176ux4ed6ux91cfux5316ux6280ux672fux6765ux5efaux7acbux548cux5206ux6790ux8fc7ux7a0bux6027ux80fdux57faux7ebfux5e76ux4fddux6301ux66f4ux65b0}}

\begin{itemize}
\tightlist
\item
  项目形成基线和标杆

  \begin{itemize}
  \tightlist
  \item
    当项目团队每轮迭代都有统计，便可以用缺陷数除以迭代的功能点规模得出缺陷率，同样也可以用返工工作量除以功能点数，在迭代里面的形成返工率。每个迭代就统计收集，并画在白纸上就形成缺陷率和返工率的每个迭代趋势图，如果发现没有太大的波动，比较稳定就可以形成项目组的基线，作为后面迭代的参考。
  \end{itemize}
\end{itemize}

\hypertarget{mpm-4.4-ux4f7fux7528ux7edfux8ba1ux4e0eux5176ux4ed6ux91cfux5316ux6280ux672fux6765ux5efaux7acbux548cux5206ux6790ux8fc7ux7a0bux6027ux80fdux6a21ux578bux5e76ux4fddux6301ux66f4ux65b0}{%
\subsubsection{MPM 4.4
使用统计与其他量化技术来建立和分析过程性能模型并保持更新}\label{mpm-4.4-ux4f7fux7528ux7edfux8ba1ux4e0eux5176ux4ed6ux91cfux5316ux6280ux672fux6765ux5efaux7acbux548cux5206ux6790ux8fc7ux7a0bux6027ux80fdux6a21ux578bux5e76ux4fddux6301ux66f4ux65b0}}

其他必需信息 Additional Required Information

\begin{description}
\tightlist
\item[]
PPM 过程性能模型：
\end{description}

\begin{itemize}
\tightlist
\item
  是根据历史过程性能数据（如过程性能基线中包含的数据）开发的

  \begin{itemize}
  \tightlist
  \item
    依据每一个过程需求评审，开发，等系统导出数据，依据数据中的缺陷建立的模型
  \end{itemize}
\item
  以可度量的属性值和术语描述、描绘变动或对其进行建模

  \begin{itemize}
  \tightlist
  \item
    有一个度量定义表，蒙特卡洛中的变量定义在度量定义表中定义;变量与预测结果都有分布
  \end{itemize}
\item
  预测中间或最终过程性能

  \begin{itemize}
  \tightlist
  \item
    蒙特卡洛可以预测出QPPO范围与过程（如需求评审）缺陷范围
  \end{itemize}
\item
  估算预测结果的预期范围和变动

  \begin{itemize}
  \tightlist
  \item
    蒙特卡洛预测出中间或最终过程QPPO的预期范围和变动
  \end{itemize}
\item
  包含至少一个表示与子过程相关的可控输入的可度量属性

  \begin{itemize}
  \tightlist
  \item
    可控因素：各阶段的评审方法（如，用走查 或 审查）
  \end{itemize}
\end{itemize}

\hypertarget{mpm-4.5-ux4f7fux7528ux7edfux8ba1ux4e0eux5176ux4ed6ux91cfux5316ux6280ux672fux6765ux786eux5b9aux6216ux9884ux6d4bux8d28ux91cfux4e0eux8fc7ux7a0bux6027ux80fdux76eeux6807ux7684ux5b9eux73b0ux60c5ux51b5}{%
\subsubsection{MPM 4.5
使用统计与其他量化技术来确定或预测质量与过程性能目标的实现情况}\label{mpm-4.5-ux4f7fux7528ux7edfux8ba1ux4e0eux5176ux4ed6ux91cfux5316ux6280ux672fux6765ux786eux5b9aux6216ux9884ux6d4bux8d28ux91cfux4e0eux8fc7ux7a0bux6027ux80fdux76eeux6807ux7684ux5b9eux73b0ux60c5ux51b5}}

\begin{itemize}
\tightlist
\item
  分析选定过程的变动和稳定性并解决缺陷

  \begin{itemize}
  \tightlist
  \item
    使用了方差分析和控制图
  \item
    通过趋势图，控制图判断历史数据是否作为基线，保证模型的稳定性。详见基线PPB截图。
  \item
    对各种方法进行方差分析，识别是否有显著区分
  \end{itemize}
\item
  实施必要的行动来解决在实现质量与过程性能目标方面的缺陷

  \begin{itemize}
  \tightlist
  \item
    P4 系统缺陷密度波动范围很广， 暂时不能预测，先对团队培训过程
  \end{itemize}
\item
  使用经数据校准并已确认的过程性能模型来评估质量与过程性能目标的实现进度

  \begin{itemize}
  \tightlist
  \item
    使用过程性能模型来预测下一轮冲刺的各阶段缺陷密度范围
  \end{itemize}
\item
  识别和管理与实现质量与过程性能目标有关的风险

  \begin{itemize}
  \tightlist
  \item
    利用PPM预测下一轮达到目标范围的概率
  \end{itemize}
\item
  记录并传达分析结果、决策以及确定的行动

  \begin{itemize}
  \tightlist
  \item
    把以上记录在量化项目管理计划
  \end{itemize}
\end{itemize}

\hypertarget{car-2.1-ux9009ux62e9ux8981ux8fdbux884cux5206ux6790ux7684ux7ed3ux679c}{%
\subsubsection{CAR 2.1
选择要进行分析的结果}\label{car-2.1-ux9009ux62e9ux8981ux8fdbux884cux5206ux6790ux7684ux7ed3ux679c}}

\begin{itemize}
\tightlist
\item
  确定哪些结果进一步分析：在迭代回顾时，对各个过程的缺陷数和遗留数进行总体分析，识别、选择关键过程来进一步分析。
\end{itemize}

\hypertarget{car-3.1-ux9075ux5faaux7ec4ux7ec7ux8fc7ux7a0bux6765ux786eux5b9aux6240ux9009ux7ed3ux679cux7684ux6839ux672cux539fux56e0}{%
\subsubsection{CAR 3.1
遵循组织过程来确定所选结果的根本原因}\label{car-3.1-ux9075ux5faaux7ec4ux7ec7ux8fc7ux7a0bux6765ux786eux5b9aux6240ux9009ux7ed3ux679cux7684ux6839ux672cux539fux56e0}}

\hypertarget{car-3.2-ux63d0ux51faux5904ux7406ux5df2ux8bc6ux522bux7684ux6839ux672cux539fux56e0ux7684ux884cux52a8ux5efaux8bae}{%
\subsubsection{CAR 3.2
提出处理已识别的根本原因的行动建议}\label{car-3.2-ux63d0ux51faux5904ux7406ux5df2ux8bc6ux522bux7684ux6839ux672cux539fux56e0ux7684ux884cux52a8ux5efaux8bae}}

\begin{itemize}
\tightlist
\item
  针对根因制定对应的新方法，例如新的评审方法或者开发方法，也估计这方法在缺陷排除或者缺陷数的量，加到模型中，利用模型的最优选择，比较这个新的方法可以比预计提升多少？
\end{itemize}

\hypertarget{car-3.3-ux5b9eux65bdux9009ux5b9aux7684ux884cux52a8ux5efaux8bae}{%
\subsubsection{CAR 3.3
实施选定的行动建议}\label{car-3.3-ux5b9eux65bdux9009ux5b9aux7684ux884cux52a8ux5efaux8bae}}

\hypertarget{car-3.4-ux8bb0ux5f55ux6839ux672cux539fux56e0ux5206ux6790ux548cux89e3ux51b3ux7684ux6570ux636e}{%
\subsubsection{CAR 3.4
记录根本原因分析和解决的数据}\label{car-3.4-ux8bb0ux5f55ux6839ux672cux539fux56e0ux5206ux6790ux548cux89e3ux51b3ux7684ux6570ux636e}}

\begin{itemize}
\tightlist
\item
  把这些都记录在A3报告里，作为公司的参考。
  公司级也可以从这些A3报告和相关数据判断是否可以扩大给其他项目组参考、使用。
\end{itemize}

\hypertarget{car-4.1-ux4f7fux7528ux7edfux8ba1ux7684ux4e0eux5176ux4ed6ux91cfux5316ux7684ux6280ux672fux5bf9ux9009ux5b9aux7ed3ux679cux8fdbux884cux6839ux672cux539fux56e0ux5206ux6790}{%
\subsubsection{CAR 4.1
使用统计的与其他量化的技术对选定结果进行根本原因分析}\label{car-4.1-ux4f7fux7528ux7edfux8ba1ux7684ux4e0eux5176ux4ed6ux91cfux5316ux7684ux6280ux672fux5bf9ux9009ux5b9aux7ed3ux679cux8fdbux884cux6839ux672cux539fux56e0ux5206ux6790}}

\begin{itemize}
\tightlist
\item
  利用参考以往基线，包括缺陷排除公司的基线，项目自己的以往迭代的基线进行比较，或借用从以往项目数据建立的公司的预测模型，帮助团队找出根因。
\end{itemize}

\hypertarget{car-4.2-ux4f7fux7528ux7edfux8ba1ux7684ux4e0eux5176ux4ed6ux91cfux5316ux7684ux6280ux672fux8bc4ux4ef7ux5b9eux65bdux7684ux884cux52a8ux5bf9ux8fc7ux7a0bux6027ux80fdux7684ux5f71ux54cd}{%
\subsubsection{CAR 4.2
使用统计的与其他量化的技术评价实施的行动对过程性能的影响}\label{car-4.2-ux4f7fux7528ux7edfux8ba1ux7684ux4e0eux5176ux4ed6ux91cfux5316ux7684ux6280ux672fux8bc4ux4ef7ux5b9eux65bdux7684ux884cux52a8ux5bf9ux8fc7ux7a0bux6027ux80fdux7684ux5f71ux54cd}}

\begin{itemize}
\tightlist
\item
  在下一个迭代执行新方法，再下一轮回顾时看效果是否有提升？比模型预计的多少？并更莫地卡罗预测模型，
\end{itemize}

\hypertarget{ux7ec4ux7ec7ux7ea7}{%
\subsection{组织级}\label{ux7ec4ux7ec7ux7ea7}}

\hypertarget{pcm-2.2-ux5236ux5b9aux4fddux6301ux66f4ux65b0ux5e76ux9075ux5faaux9009ux5b9aux7684ux8fc7ux7a0bux6539ux8fdbux4e0eux5b9eux65bdux8ba1ux5212}{%
\subsubsection{PCM 2.2
制定、保持更新并遵循选定的过程改进与实施计划}\label{pcm-2.2-ux5236ux5b9aux4fddux6301ux66f4ux65b0ux5e76ux9075ux5faaux9009ux5b9aux7684ux8fc7ux7a0bux6539ux8fdbux4e0eux5b9eux65bdux8ba1ux5212}}

\begin{itemize}
\tightlist
\item
  过程改进计划

  \begin{itemize}
  \tightlist
  \item
    与干系人开会讨论后发改进项目概览页（POS）
  \item
    用迭代原型法制定过程改进计划，与试点项目每轮迭代改进同步
  \end{itemize}
\end{itemize}

\hypertarget{pcm-3.1-ux5236ux5b9aux6301ux7eedux66f4ux65b0ux5e76ux4f7fux7528ux53efux8ffdux6eafux5230ux4e1aux52a1ux76eeux6807ux7684ux8fc7ux7a0bux6539ux8fdbux76eeux6807}{%
\subsubsection{PCM 3.1
制定、持续更新并使用可追溯到业务目标的过程改进目标}\label{pcm-3.1-ux5236ux5b9aux6301ux7eedux66f4ux65b0ux5e76ux4f7fux7528ux53efux8ffdux6eafux5230ux4e1aux52a1ux76eeux6807ux7684ux8fc7ux7a0bux6539ux8fdbux76eeux6807}}

\begin{itemize}
\tightlist
\item
  参考业务目标，制定组织级过程改进目标

  \begin{itemize}
  \tightlist
  \item
    开始用减小缺陷返工工作量，预估能为公司节省多少，说服高层支持等，把过程改进正式立项。{[}第八章{]}
  \item
    当项目试点开始后，开始有实际的迭代数据组织级的过程改进组便应按实际的试点项目结果更新目标。{[}第七章{]}
  \end{itemize}
\end{itemize}

\hypertarget{pcm-3.2-ux8bc6ux522bux6700ux6709ux52a9ux4e8eux5b9eux73b0ux4e1aux52a1ux76eeux6807ux7684ux8fc7ux7a0b}{%
\subsubsection{PCM 3.2
识别最有助于实现业务目标的过程}\label{pcm-3.2-ux8bc6ux522bux6700ux6709ux52a9ux4e8eux5b9eux73b0ux4e1aux52a1ux76eeux6807ux7684ux8fc7ux7a0b}}

\begin{itemize}
\tightlist
\item
  利用二八原则时找出应针对哪方面做根因分析和对应的纠正措施。

  \begin{itemize}
  \tightlist
  \item
    可以从子过程、人员等去利用二八原则，找出关键的改进点，应先从试点项目入手{[}十一章{]}
  \end{itemize}
\end{itemize}

\hypertarget{pcm-3.4-ux63a2ux7d22ux548cux8bc4ux4f30ux6f5cux5728ux7684ux65b0ux8fc7ux7a0bux6280ux672fux65b9ux6cd5ux548cux5de5ux5177ux6765ux8bc6ux522bux6539ux8fdbux673aux4f1a}{%
\subsubsection{PCM 3.4
探索和评估潜在的新过程、技术、方法和工具来识别改进机会}\label{pcm-3.4-ux63a2ux7d22ux548cux8bc4ux4f30ux6f5cux5728ux7684ux65b0ux8fc7ux7a0bux6280ux672fux65b9ux6cd5ux548cux5de5ux5177ux6765ux8bc6ux522bux6539ux8fdbux673aux4f1a}}

\begin{itemize}
\tightlist
\item
  挑选试点项目和试点后的数据分析

  \begin{itemize}
  \tightlist
  \item
    在公司立项后便应做试点，在试点的每轮迭代分析缺陷数据和返工数据，形成A3报告。这些在前面上面CAR都介绍过了，组机级的改进组应收集各个试点项目的报告跟高层汇报看跟本来的预期是相差多少，是否需要更新计划{[}第七章{]}
  \end{itemize}
\end{itemize}


